\DocumentMetadata{
	tagging=on,
	tagging-setup={math/setup=mathml-SE},
	pdfstandard=ua-2,
	pdfversion=2.0,
	lang=en
}

\documentclass[hidelinks, 11 pt]{article}

\usepackage[a4paper, top=25.4mm, bottom=25.4mm, left=19.1mm, right=19.1mm, footskip=8mm]{geometry}
\usepackage{mathtools}
\usepackage{pgfplots}
\usepackage{hyperref}
\usepackage[list=true, listformat=simple]{subcaption}
\usepackage{csvsimple}
\usepackage{booktabs}
\usepackage{longtable}
\usepackage{biblatex}
\usepackage{fancyhdr}

\addbibresource{bibliography.bib}

\pgfplotsset{compat=1.17}

\newcommand{\myhline}{\hline\noalign{\vskip 2pt}} % Space after hline
\renewcommand{\arraystretch}{1.2} % Increase spacing between rows of tables
\hypersetup{
	pdfdisplaydoctitle=true
}


% Dark mode

\definecolor{myblack}{HTML}{1E2021}
\definecolor{mywhite}{HTML}{E4E5ED}
\newcommand{\AnnotationColour}{myblack}

\newcommand{\EnableDifferenceTikz}{
	\AtBeginDocument{
		\IfPackageLoadedTF{tikz}{
			% Annotation colour for dark mode
			\renewcommand{\AnnotationColour}{mywhite}
			
			\tikzset{
				every node/.append style={
					blend mode=difference
				}
			}
			
			\pgfplotsset{
				every axis/.append style={
					blend mode=normal
				},
				legend style={
					blend mode=normal,
					fill=myblack,
					draw=mywhite
				}
			}
		}{}
	}
}

\newcommand{\StyleNormal}{
	\hypersetup{
		colorlinks=true,
		urlcolor=blue,
		linkcolor=blue,
		citecolor=blue,
		filecolor=blue
	}
}

\newcommand{\StyleDark}{
	\pagecolor{myblack}
	\color{mywhite}
	
	\colorlet{mylinkblue}{blue!10}
	
	\hypersetup{
		colorlinks=true,
		urlcolor=mylinkblue,
		linkcolor=mylinkblue,
		citecolor=mylinkblue,
		filecolor=mylinkblue,
	}
	
	\EnableDifferenceTikz
}

%\StyleNormal
\StyleDark

\makeatletter
\setlength{\@fptop}{0pt plus 1fil}
\setlength{\@fpsep}{0pt plus 1fil}
\setlength{\@fpbot}{0pt plus 1fil}
\makeatother

% Defining marks
\renewcommand{\sectionmark}[1]{\gdef\leftmark{#1}}

\fancypagestyle{firstpage}{
	\fancyhf{}
	\fancyfoot[C]{\thepage}
	\renewcommand{\headrulewidth}{0pt}
}

\fancypagestyle{fancy}{
	\fancyhf{}
	\fancyhead[L]{\nouppercase{\leftmark}}
	\fancyhead[R]{\nouppercase{Monopoly Run}}
	\fancyfoot[C]{\thepage}
}

\pagestyle{fancy}

% Line spacing between paragraphs

\usepackage[skip=0.5\baselineskip]{parskip}

\makeatletter
\renewcommand\section{%
	\@startsection{section}{1}{\z@}%
	{-2.5ex plus -1ex minus -.2ex}%
	{1.5ex plus .2ex}%
	{\normalfont\Large\bfseries}}

\renewcommand\subsection{%
	\@startsection{subsection}{2}{\z@}%
	{-2.0ex plus -0.8ex minus -.2ex}%
	{1.0ex plus .2ex}%
	{\normalfont\large\bfseries}}

\renewcommand\subsubsection{%
	\@startsection{subsubsection}{3}{\z@}%
	{-1.5ex plus -0.6ex minus -.2ex}%
	{0.8ex plus .2ex}%
	{\normalfont\normalsize\bfseries}}
\makeatother

\newcommand{\RouteFigure}[3]{%
	\begin{subfigure}{#3\textwidth}
		\centering
		\begin{tikzpicture}
			\begin{axis}[
				width=\linewidth,
				axis equal,
				hide axis,
				clip=false,
				]
				
				% Route
				\addplot[
				purple,
				thick,
				mark=none,
				] table[
				x index=0,
				y index=1,
				col sep=space,
				header=false,
				] {../Output/PlottingEdges/#1.csv};
				
				% Places
				\addplot[
				purple,
				only marks,
				mark=*,
				mark size=2pt,
				nodes near coords,
				point meta=explicit symbolic,
				every node near coord/.append style={
					text=\AnnotationColour,
					font=\footnotesize,
					anchor=center,
				},
				] table[
				x=X,
				y=Y,
				meta=Place,
				col sep=comma,
				] {../Output/PlottingPlaces/#1.csv};
				
				% Labels
				%				\addplot[
				%				only marks,
				%				mark=none,
				%				nodes near coords,
				%				point meta=explicit symbolic,
				%				every node near coord/.append style={
					%					text=\AnnotationColour,
					%					font=\footnotesize,
					%					anchor=center,
					%				},
				%				] table[
				%				x=LabelX,
				%				y=LabelY,
				%				meta=Place,
				%				col sep=comma,
				%				] {../Output/PlottingPlaces/#1.csv};
				%				
			\end{axis}
		\end{tikzpicture}%
		\caption{#2 m/s}
		\label{fig:Solution_#1}
	\end{subfigure}
}

\ExplSyntaxOn
\renewcommand{\thesubfigure}{\int_to_alph:n{\value{subfigure}}}
\ExplSyntaxOff